\subsection{训练方案}
\label{sec:training}

网络使用 AdamW 优化器~\cite{Loshchilov2019}进行训练，初始学习率为$1 \times 10^{-3}$，解耦权重衰减设为$1 \times 10^{-4}$，小批量大小为$32$。学习率遵循余弦退火调度，从$1 \times 10^{-3}$至$1 \times 10^{-4}$覆盖整个训练过程。Dropout 概率$p = 0.2$应用于融合层；L2 正则化（权重衰减）进一步约束过拟合。禁用批归一化以保留声学信号的频率特异性。

\textbf{数据增强。}训练期间，每个输入信号在样本基础上进行随机增强：随机时间平移为信号长度的$\pm 12.5\%$，相对于干净信号的信噪比(SNR)从$[-6, +6]$~dB 均匀采样的加性高斯噪声注入，以及从$[0.8, 1.2]$均匀采样的幅值缩放。增强独立应用于有限平均输入和高平均目标，以保持配对完整性。

\textbf{验证和早停。}训练集由完整铝试样成对数据的 80\% 组成；保留 20\% 作为验证集用于监控收敛。
当验证重构损失在连续 15 个 epoch 无改善时训练终止，恢复最佳验证模型权重。
在使用的硬件配置上（单块 NVIDIA RTX 4090，24~GB 显存），完整训练运行约在 180 个 epoch 内完成。

\textbf{模型规模。}PMDF-Net 包含约$4.2 \times 10^6$个可训练参数。通过两个分支和融合模块的前向传播对于长度$L = 1024$的输入信号产生约$1.8 \times 10^9$次浮点运算（FLOPs）。在 RTX 4090 上的推理延迟为每信号$2.3 \pm 0.1$~ms，满足实际检测吞吐量的实时性要求。